\thispagestyle{empty}

\begin{center}
  \textbf{UNIVERSIDADE FEDERAL DO PIAUÍ (UFPI)}\\
  \textbf{CENTRO DE CIÊNCIAS DA NATUREZA (CCN)}\\
  \textbf{DEPARTAMENTO DE COMPUTAÇÃO (DC)}\\[0.3em]
  Disciplina: Tópicos em Inteligência Artificial\\
  Professor: Raimundo Santos Moura \quad--\quad Período: 2026.1
\end{center}

\vspace{3.5cm}

\begin{center}
  {\Large \textbf{Pré-Treino e Pós-Treino de Modelos de Linguagem para o Domínio Público Piauiense}}\\[0.8em]
  {\large Relatório das Questões 01 (Pré-Treino Continuado), 02 (SFT) e 03 (LoRA/QLoRA)}
\end{center}

\vspace{2.5cm}

\begin{center}
\begin{tabularx}{0.8\textwidth}{|>{\RaggedRight\arraybackslash}X|}
  \hline
  \textbf{Modelo base:} Qwen2.5-0.5B (494M parâmetros)\\
  \hline
  \textbf{Datasets:} DOMPI-2025 (Diário Oficial dos Municípios do Piauí) e docentesDC (materiais dos docentes do DC/UFPI)\\
  \hline
  \textbf{Ambiente de treino:} Kaggle (1 GPU NVIDIA T4, 16~GB)\\
  \hline
\end{tabularx}
\end{center}

\vspace{3cm}

\begin{center}
  \textbf{Autores}\\[0.3em]
  Allycia \quad\textbar\quad Heverton \quad\textbar\quad Izaías \quad\textbar\quad Oscar
\end{center}

\vfill

\begin{center}
  Teresina -- PI\\
  2026
\end{center}
